% ==========================================================================
% Kapitel 1: Logik
% ==========================================================================

\section{Kapitel 1: Logik}

\vorlesungsdatum{26.10.2026}{handschrift/26102026.pdf}

In diesem Kapitel besprechen wir kurz und prägnant die aussagen- und prädikatenlogischen Grundlagen, die nötig sind, um Mathematik formal und widerspruchsfrei zu betreiben. Die Logik dient uns hierbei vor allem als universelle Sprache zur Formulierung von Definitionen, Sätzen und Beweisen.

% --------------------------------------------------------------------------
% Aussagen
% --------------------------------------------------------------------------
\subsection{Aussagen}

Ausgangspunkt der mathematischen Logik ist der Begriff der Aussage:

\begin{definition}[Aussage]
Eine \textbf{Aussage} ist ein sprachlicher, grammatikalisch korrekter Ausdruck (Satz), welcher eindeutig \textbf{wahr} ($1$ bzw.\ $w$) oder eindeutig \textbf{falsch} ($0$ bzw.\ $f$) ist.
\end{definition}

Die klassische Logik beruht auf dem fundamentalen \textbf{Prinzip der Zweiwertigkeit}, welches sich in zwei Kernprinzipien gliedert:
\begin{enumerate}
    \item \textbf{Prinzip vom ausgeschlossenen Widerspruch:} Eine Aussage kann nicht zugleich wahr und falsch sein.
    \item \textbf{Satz vom ausgeschlossenen Dritten (\textit{Tertium non datur}):} Eine Aussage hat stets genau einen der beiden Wahrheitswerte \glqq wahr\grqq\ oder \glqq falsch\grqq; Zwischenwerte wie \glqq vielleicht\grqq\ existieren nicht.
\end{enumerate}

\begin{bemerkung}[Abgrenzung zu Nicht-Aussagen]
Nicht jeder sprachlich wohlgeformte Satz stellt eine mathematische Aussage dar:
\begin{itemize}
    \item Fragen (\glqq Wie spät ist es?\grqq)
    \item Meinungsäußerungen und Werturteile (\glqq Dieses Bild ist schön.\grqq)
    \item Aufforderungen und Befehle (\glqq Schließe bitte die Tür!\grqq)
    \item Terme ohne Relationszeichen: $9 + 3$ ist keine Aussage (sondern ein Term); die Gleichung $9 + 3 = 12$ hingegen ist eine wahre Aussage.
    \item Paradoxien wie das Lügner-Paradoxon (\glqq Dieser Satz ist falsch.\grqq).
\end{itemize}
\end{bemerkung}

\begin{beispiel}[Wahrheitswert und Prüfbarkeit]
Ob ein Ausdruck eine Aussage ist, hängt nicht davon ab, ob der Wahrheitswert unmittelbar auf einen Blick ersichtlich ist:
\begin{itemize}
    \item \glqq $31\,274\,567$ ist eine Primzahl.\grqq\ ist eine Aussage (sie ist objektiv wahr oder falsch, auch wenn die Prüfung Rechenaufwand erfordert).
    \item \glqq Die Erde ist rund.\grqq\ $\rightarrow$ wahre Aussage.
    \item \glqq Dieser Satz besteht aus sechs Wörtern.\grqq\ $\rightarrow$ wahre Aussage (selbstreferenziell, aber eindeutig entscheidbar).
\end{itemize}
\end{beispiel}

\subsection*{Das Extensionalitätsprinzip}
\textbf{Extensionalitätsprinzip (Wahrheitsfunktionalität):} Der Wahrheitswert einer zusammengesetzten Aussage hängt ausschließlich von den Wahrheitswerten ihrer Teilkomponenten ab -- nicht von inhaltlichen, zeitlichen oder kausalen Zusammenhängen.

\begin{bemerkung}[Unterschied zur Umgangssprache]
In der Umgangssprache transportiert die Satzreihenfolge bei \glqq und\grqq\ häufig Kausalität oder zeitliche Abfolge:
\begin{enumerate}
    \item \glqq Sie verließ das Land und wurde reich.\grqq
    \item \glqq Sie wurde reich und verließ das Land.\grqq
\end{enumerate}
In der formalen Logik ist die Verknüpfung $A \land B$ rein wahrheitswertbasiert und kommutativ ($A \land B \equiv B \land A$); Zeit- und Kausalbezüge spielen keine Rolle.
\end{bemerkung}

% --------------------------------------------------------------------------
% Verknuepfungen und Wahrheitstafeln
% --------------------------------------------------------------------------
\subsection{Logische Verknüpfungen und Wahrheitstafeln}

Seien $A$ und $B$ Aussagen. Neue, zusammengesetzte Aussagen werden über \textbf{Wahrheitstafeln} eindeutig festgelegt:

\begin{definition}[Logische Verknüpfungen]
Seien $A$ und $B$ Aussagen. Wir definieren:
\begin{enumerate}
    \item \textbf{Konjunktion ($A \land B$, \glqq $A$ und $B$\grqq):} Nur dann wahr ($1$), wenn sowohl $A$ als auch $B$ wahr sind.
    \item \textbf{Disjunktion ($A \lor B$, \glqq $A$ oder $B$\grqq):} Inklusives Oder; wahr, sobald mindestens eine der Aussagen $A$ oder $B$ wahr ist.
    \item \textbf{Antivalenz ($A \xor B$, \glqq Entweder $A$ oder $B$\grqq, XOR):} Genau dann wahr, wenn $A$ und $B$ unterschiedliche Wahrheitswerte haben.
    \item \textbf{Negation ($\neg A$ bzw.\ $\bar{A}$, \glqq Nicht $A$\grqq):} Kehrt den Wahrheitswert der Aussage um.
    \item \textbf{Implikation ($A \Rightarrow B$, \glqq Wenn $A$, dann $B$\grqq):} Nur dann falsch, wenn $A$ wahr ($1$) und $B$ falsch ($0$) ist. Ist $A$ falsch, ist die Implikation stets wahr (\emph{Ex falso quodlibet}). $A$ heißt hinreichend für $B$; $B$ heißt notwendig für $A$.
    \item \textbf{Äquivalenz ($A \Leftrightarrow B$, \glqq $A$ genau dann, wenn $B$\grqq):} Wahr, wenn $A$ und $B$ denselben Wahrheitswert besitzen.
\end{enumerate}
\end{definition}

\begin{center}
\renewcommand{\arraystretch}{1.15}
\begin{tabular}{ccccccc}
    \toprule
    $A$ & $B$ & $A \land B$ & $A \lor B$ & $A \xor B$ & $A \Rightarrow B$ & $A \Leftrightarrow B$ \\
    \midrule
    0 & 0 & 0 & 0 & 0 & 1 & 1 \\
    0 & 1 & 0 & 1 & 1 & 1 & 0 \\
    1 & 0 & 0 & 1 & 1 & 0 & 0 \\
    1 & 1 & 1 & 1 & 0 & 1 & 1 \\
    \bottomrule
\end{tabular}
\qquad
\begin{tabular}{cc}
    \toprule
    $A$ & $\neg A$ \\
    \midrule
    0 & 1 \\
    1 & 0 \\
    \bottomrule
\end{tabular}
\end{center}

\begin{bemerkung}[Ex falso quodlibet]
Die Belegung von $A \Rightarrow B$ bei falscher Prämisse ($A = 0$) lässt sich an Beispielen motivieren:
\begin{itemize}
    \item \glqq Wenn ich heute im Lotto gewinne, mache ich eine Weltreise.\grqq\newline
    Gewinne ich nicht, habe ich mein Versprechen nicht gebrochen -- die Aussage bleibt wahr.
    \item Mathematisch: Die Aussagen \glqq Wenn $2+2=42$, dann ist $2+2=0$\grqq\ und \glqq Wenn $2+2=42$, dann ist die Erde eine Scheibe\grqq\ sind formal beide \textbf{wahr}, da die Voraussetzung falsch ist.
\end{itemize}
\end{bemerkung}

\subsection*{Operatorrangfolge}
Um Klammern zu sparen, gilt folgende Bindungshierarchie:
\begin{center}
1.~Negation $\neg$ \quad $>$ \quad 2.~Konjunktion / Disjunktion $\land, \lor$ \quad $>$ \quad 3.~Implikation / Äquivalenz $\Rightarrow, \Leftrightarrow$
\end{center}
\emph{Beispiel:} $\neg A \lor B$ bedeutet stets $(\neg A) \lor B$ und nicht $\neg(A \lor B)$.

\begin{definition}[Logische Äquivalenz $\equiv$]
Zwei zusammengesetzte Aussagen $A$ und $B$ heißen \textbf{logisch äquivalent} (in Zeichen $A \equiv B$), wenn sie bei jeder Belegung der Teilkomponenten denselben Wahrheitswert liefern.
\end{definition}

\begin{satz}[Auflösung der Implikation]
Für beliebige Aussagen $A$ und $B$ gilt:
\[
(A \Rightarrow B) \;\equiv\; (\neg A \lor B).
\]
\end{satz}

\begin{proof}
Die Wahrheitstafel zeigt zeilenweise identische Ergebnisse:
\begin{center}
\renewcommand{\arraystretch}{1.15}
\begin{tabular}{ccccc}
    \toprule
    $A$ & $B$ & $\neg A$ & $A \Rightarrow B$ & $\neg A \lor B$ \\
    \midrule
    0 & 0 & 1 & 1 & 1 \\
    0 & 1 & 1 & 1 & 1 \\
    1 & 0 & 0 & 0 & 0 \\
    1 & 1 & 0 & 1 & 1 \\
    \bottomrule
\end{tabular}
\end{center}
Da die Spalten für $A \Rightarrow B$ und $\neg A \lor B$ übereinstimmen, ist die logische Äquivalenz bewiesen.
\end{proof}

% --------------------------------------------------------------------------
% Tautologien und Schlussregeln
% --------------------------------------------------------------------------
\subsection{Tautologien und Schlussregeln}

\begin{definition}[Tautologie und Widerspruch]
Eine zusammengesetzte Aussage heißt \textbf{Tautologie}, wenn sie unabhängig von den Wahrheitswerten ihrer Teilaussagen stets den Wert \glqq wahr\grqq\ ($1$) annimmt.\newline
Sie heißt \textbf{Widerspruch} (Symbol $\bot$), wenn sie stets den Wert \glqq falsch\grqq\ ($0$) annimmt.
\end{definition}

Es gilt also: $A \equiv B$ genau dann, wenn $A \Leftrightarrow B$ eine Tautologie ist.

\begin{satz}[De Morgansche Regeln]
Für beliebige Aussagen $A$ und $B$ gelten:
\begin{align*}
    \neg(A \land B) &\Longleftrightarrow (\neg A \lor \neg B), \\
    \neg(A \lor B) &\Longleftrightarrow (\neg A \land \neg B).
\end{align*}
\end{satz}

\begin{proof}
Beweis der ersten Regel exemplarisch mittels Wahrheitstafel:
\begin{center}
\renewcommand{\arraystretch}{1.15}
\begin{tabular}{cccccccc}
    \toprule
    $A$ & $B$ & $\neg A$ & $\neg B$ & $A \land B$ & $\neg(A \land B)$ & $\neg A \lor \neg B$ & $\neg(A \land B) \Leftrightarrow (\neg A \lor \neg B)$ \\
    \midrule
    0 & 0 & 1 & 1 & 0 & 1 & 1 & 1 \\
    0 & 1 & 1 & 0 & 0 & 1 & 1 & 1 \\
    1 & 0 & 0 & 1 & 0 & 1 & 1 & 1 \\
    1 & 1 & 0 & 0 & 1 & 0 & 0 & 1 \\
    \bottomrule
\end{tabular}
\end{center}
Da die letzte Spalte ausschließlich den Wert 1 aufweist, handelt es sich um eine Tautologie.
\end{proof}

\begin{satz}[Klassische Schlussregeln]
Es gelten die folgenden fundamentalen Tautologien:
\begin{enumerate}
    \item \textbf{Doppelte Negation:} $A \Longleftrightarrow \neg(\neg A)$
    \item \textbf{Modus Ponens:} $(A \land (A \Rightarrow B)) \Longrightarrow B$
    \item \textbf{Gesetz des Syllogismus (Transitivität):} $((A \Rightarrow B) \land (B \Rightarrow C)) \Longrightarrow (A \Rightarrow C)$
    \item \textbf{Kontraposition (Indirekter Beweis 1):} $(A \Rightarrow B) \Longleftrightarrow (\neg B \Rightarrow \neg A)$
    \item \textbf{Widerspruchsbeweis (Indirekter Beweis 2):} $A \Longleftrightarrow (\neg A \Rightarrow \bot)$
    \item \textbf{Äquivalenzbeweis:} $(A \Leftrightarrow B) \Longleftrightarrow ((A \Rightarrow B) \land (B \Rightarrow A))$
    \item \textbf{Distributivgesetze:}
    \begin{align*}
        A \land (B \lor C) &\Longleftrightarrow (A \land B) \lor (A \land C), \\
        A \lor (B \land C) &\Longleftrightarrow (A \lor B) \land (A \lor C).
    \end{align*}
\end{enumerate}
\end{satz}

\begin{proof}[Beweis der Kontraposition ohne Wahrheitstafel]
Unter wiederholter Anwendung bereits bewiesener Äquivalenzen:
\[
(A \Rightarrow B) \;\Longleftrightarrow\; (\neg A \lor B) \;\Longleftrightarrow\; (\neg A \lor \neg\neg B) \;\Longleftrightarrow\; (\neg(\neg B) \lor \neg A) \;\Longleftrightarrow\; (\neg B \Rightarrow \neg A).
\]
\end{proof}

\subsection*{Klassische Anwendung: Irrationalität von \texorpdfstring{$\sqrt{2}$}{sqrt(2)}}

\begin{satz}[Irrationalität von $\sqrt{2}$]
Die Zahl $\sqrt{2}$ ist irrational ($\sqrt{2} \notin \Q$).
\end{satz}

\begin{proof}
Wir führen einen indirekten Beweis (Widerspruchsbeweis): Angenommen, $\sqrt{2}$ sei rational. Dann existieren teilerfremde natürliche Zahlen $p, q \in \N_{>0}$ mit:
\[
\sqrt{2} = \frac{p}{q}.
\]
Quadrieren und Multiplikation mit $q^2$ ergibt:
\[
2q^2 = p^2.
\]
Hieraus folgt, dass $p^2$ gerade ist. Da das Quadrat einer ungeraden Zahl stets ungerade ist, muss auch $p$ gerade sein. Es existiert also ein $r \in \N$ mit $p = 2r$. Einsetzen liefert:
\[
2q^2 = (2r)^2 = 4r^2 \implies q^2 = 2r^2.
\]
Dasselbe Argument zeigt, dass auch $q$ gerade sein muss. 
Damit sind sowohl $p$ als auch $q$ durch 2 teilbar, was der Teilerfremdheit von $p$ und $q$ widerspricht. Folglich ist die Annahme falsch und $\sqrt{2}$ ist irrational.
\end{proof}

% --------------------------------------------------------------------------
% Aussageformen und Quantoren
% --------------------------------------------------------------------------
\subsection{Aussageformen und Quantoren}

\begin{definition}[Aussageform]
Eine \textbf{einstellige Aussageform} mit Einsetzungsklasse $\Omega$ ist ein sprachlicher, grammatikalisch korrekter Ausdruck $H(\cdot)$ mit einer Leerstelle, sodass für jedes $x \in \Omega$ die Einsetzung $H(x)$ eine Aussage ist.
\end{definition}

\begin{beispiel}
\begin{itemize}
    \item Sei $\Omega = \R$ und $H(x) :\Longleftrightarrow x \in \Q$. Dann ist $H(42)$ wahr, $H(\sqrt{2})$ jedoch falsch.
    \item Sei $\Omega = \R$ und $H(x) :\Longleftrightarrow x^2 \ge 0$. Dies ist für alle $x \in \R$ wahr.
    \item Sei $\Omega = \N$ und $H(n) :\Longleftrightarrow 2n+4 \text{ ist Summe zweier Primzahlen}$ (Goldbachsche Vermutung).
\end{itemize}
\end{beispiel}

\begin{definition}[All- und Existenzquantor]
Ist $H(\cdot)$ eine Aussageform über $\Omega$, so bilden wir:
\begin{itemize}
    \item $\forall x \in \Omega : H(x)$ \quad (\glqq Für alle $x \in \Omega$ gilt $H(x)$\grqq, Allquantor)
    \item $\exists x \in \Omega : H(x)$ \quad (\glqq Es gibt ein $x \in \Omega$, sodass $H(x)$ gilt\grqq, Existenzquantor)
    \item $\exists! x \in \Omega : H(x)$ \quad (\glqq Es gibt genau ein $x \in \Omega$, sodass $H(x)$ gilt\grqq)
\end{itemize}
\end{definition}

\begin{satz}[Negation von Quantoren]
Es gelten die Äquivalenzen:
\begin{align*}
    \neg(\forall x \in \Omega : H(x)) &\Longleftrightarrow \exists x \in \Omega : \neg H(x), \\
    \neg(\exists x \in \Omega : H(x)) &\Longleftrightarrow \forall x \in \Omega : \neg H(x).
\end{align*}
\end{satz}

\begin{bemerkung}[Quantoren auf Teilklassen]
Für eine Unterklasse $\Omega' \subset \Omega$ gilt:
\begin{align*}
    (\exists x \in \Omega' : H(x)) &\Longleftrightarrow \exists x \in \Omega : (x \in \Omega' \land H(x)), \\
    (\forall x \in \Omega' : H(x)) &\Longleftrightarrow \forall x \in \Omega : (x \in \Omega' \Longrightarrow H(x)).
\end{align*}
\emph{Wichtig:} Bei der Allaussage steht eine Implikation. Daher ist ein Satz wie \glqq Alle runden Dreiecke sind blau\grqq\ im mathematischen Sinne formal wahr, da die Prämisse (rundes Dreieck) niemals erfüllt ist.
\end{bemerkung}

\subsection*{Reihenfolge bei mehreren Quantoren}
Treffen unterschiedliche Quantoren aufeinander, ist die Reihenfolge entscheidend:
\begin{beispiel}
Sei $\Omega$ die Menge aller verheirateten Personen und
\[
H(x, y) :\Longleftrightarrow x \text{ ist mit } y \text{ verheiratet}.
\]
\begin{itemize}
    \item $\forall x \in \Omega \; \exists y \in \Omega : H(x, y)$:\enspace \glqq Jeder hat einen Ehepartner\grqq\ $\rightarrow$ \textbf{wahr}.
    \item $\exists y \in \Omega \; \forall x \in \Omega : H(x, y)$:\enspace \glqq Es gibt jemanden, der mit allen verheiratet ist\grqq\ $\rightarrow$ \textbf{falsch}.
\end{itemize}
\end{beispiel}

\vspace{0.8em}
\hrule
\vspace{1em}
